% Reference entries, B (AN-11).  Law 1: os/iword/app/face_a.c block_key
% (the keys that act while Block is on; "Block on", "Delete Block? No (Yes)",
% M_FONTBLK "Attribute: 1 Size; 2 Appearance: 0"), os/iword/core/block.c
% (the block from the anchor to the cursor; iw_ed_block_attr), codes.c
% iw_ed_attr (F6: the pair of codes, the cursor between them; F6 again moves
% past the off code).  Shots: ../wordperfect51-decomp mint MINT12, MINT28.
% GATES: test-pubs-iword test-iword-codes test-iword-face test-flair-iword-block test-flair-iword-format
\Gates{test-pubs-iword test-iword-codes test-iword-face test-flair-iword-block test-flair-iword-format}
\begin{Feature}{Block}
Block marks a stretch of text so that you can act on all of it at once:
make it bold or underlined, give it a size or an appearance, move it, copy
it or delete it.

You turn Block on with the cursor at one end of the text and move the
cursor to the other end.  The text between the place where you turned Block
on and the cursor is the block.  It is shown highlighted, and \Scr{Block on}
is shown at the left of the status line while Block is on.

To block text,
\begin{WPSteps}
\item Move the cursor to the first (or the last) character of the text.
\item Press \FK{Block}{Alt-F4}.  \Scr{Block on} appears on the status line.
\item Move the cursor to the other end of the text with the cursor keys.
The highlight follows the cursor.
\item Press the key of the feature you want to use on the block: Bold
(\K{F6}), Underline (\K{F8}), Font (\K{Ctrl-F8}), Move (\K{Ctrl-F4}), or
\K{Delete}.
\end{WPSteps}
\FlairShot{block}{640}{440}{Block on: the first two words of the sentence blocked}

\Documents{wpkey}{Alt-F4}
\Documents{wpkey}{Block:Alt-F4}
\Documents{wpkey}{Block:F1}
\Documents{wpkey}{Block:Esc}
\Documents{wpkey}{Block:Alt-F3}
To turn Block off without using it, press \FK{Block}{Alt-F4} again, or
press \FK{Cancel}{F1} or \K{Esc}.  The text stays as it was.

\Notes
\Topic{Cursor Keys While Block Is On}
Every cursor key works while Block is on, including the Home combinations
(see \InRef{Cursor Movement}): \K{Home, Right} or \K{End} stretches the
block to the end of the line, \K{Home, Home, Down} to the end of the
document.  The block may run backward from where you turned Block on as
well as forward.

\Topic{Features That Use a Block}
While Block is on, InitechWord acts only on these keys: the cursor keys,
\K{Alt-F4}, \K{F1} and \K{Esc} (Block off), \K{F6} and \K{F8} (see
\InRef{Bold} and \InRef{Underline}), \K{Ctrl-F8} (see \InRef{Block,
Attributes}), \K{Ctrl-F4} (see \InRef{Move, Block}), \K{Delete} and
\K{Backspace} (see \InRef{Block, Delete}), and \K{Alt-F3} (Reveal Codes).
Typing a character does nothing while Block is on.

\Topic{Highlight}
The block is drawn in black on grey.  The cell under the cursor is drawn in
reverse video, as always, and is not part of the highlight: the block ends
just before the cursor when the cursor is past its end.

\Topic{Reveal Codes}
With Reveal Codes on, the codes area shows \Code{[Block]} where you turned
Block on.  \Code{[Block]} is not a code in the document; it goes away when
Block is turned off.

\Topic{Selecting a Sentence, Paragraph or Page}
Instead of blocking by hand, you can have InitechWord block the sentence,
paragraph or page the cursor is in (see \InRef{Move, Sentence/Paragraph/Page}).

\SeeAlso{Block, Attributes; Block, Delete; Cancel; Move, Block; Reveal Codes}
\end{Feature}

\begin{Feature}{Block, Attributes}
You can give a block any appearance or size: InitechWord puts the
attribute's on code at the start of the block and its off code at the end,
and turns Block off.

To give a block an attribute,
\begin{WPSteps}
\item Block the text (see \InRef{Block}).
\item Press \FK{Bold}{F6} to make the block bold
\Or
Press \FK{Underline}{F8} to underline it
\Or
Press \FK{Font}{Ctrl-F8}.  The status line shows
\Scr{Attribute: 1 Size; 2 Appearance: 0}.  Select \Ch{Size}{1} or
\Ch{Appearance}{2}, then the size or appearance by its number, as on the
Size and Appearance menus.
\end{WPSteps}
\Documents{wpkey}{Block:F6}
\Documents{wpkey}{Block:F8}
\Documents{wpkey}{Block:Ctrl-F8}
\Documents{wpmenu}{Block:Ctrl-F8,1=Size}
\Documents{wpmenu}{Block:Ctrl-F8,1,1=Suprscpt}
\Documents{wpmenu}{Block:Ctrl-F8,1,2=Subscpt}
\Documents{wpmenu}{Block:Ctrl-F8,1,3=Fine}
\Documents{wpmenu}{Block:Ctrl-F8,1,4=Small}
\Documents{wpmenu}{Block:Ctrl-F8,1,5=Large}
\Documents{wpmenu}{Block:Ctrl-F8,1,6=Vry Large}
\Documents{wpmenu}{Block:Ctrl-F8,1,7=Ext Large}
\Documents{wpmenu}{Block:Ctrl-F8,2=Appearance}
\Documents{wpmenu}{Block:Ctrl-F8,2,1=Bold}
\Documents{wpmenu}{Block:Ctrl-F8,2,2=Undln}
\Documents{wpmenu}{Block:Ctrl-F8,2,3=Dbl Und}
\Documents{wpmenu}{Block:Ctrl-F8,2,4=Italc}
\Documents{wpmenu}{Block:Ctrl-F8,2,5=Outln}
\Documents{wpmenu}{Block:Ctrl-F8,2,6=Shadw}
\Documents{wpmenu}{Block:Ctrl-F8,2,7=Sm Cap}
\Documents{wpmenu}{Block:Ctrl-F8,2,8=Redln}
\Documents{wpmenu}{Block:Ctrl-F8,2,9=Stkout}

For example, to make the first word of a sentence bold after you have
typed it,
\BeginReplay{ref-blockattr}
\begin{WPSteps}
\item Type \TypeRun{The report is late.}
\item Press \PressKey{Home}, \PressKey{Left} to move to the start of the line.
\item Press \Press{Block}{Alt-F4}, then press \PressKey{Right} three times.\Keys{Right Right}
\WordScreen[rows=1]{ref-blockattr-1}{The first word blocked}
\item Press \Press{Bold}{F6}.  The word is bold and Block is off.
\item Press \Press{Reveal Codes}{Alt-F3}.
\end{WPSteps}
\WordScreen[rows=13]{ref-blockattr-2}{The block made bold: the codes at its two ends}
\EndReplay

\Notes
\Topic{The Cursor}
After the attribute is given, the cursor is at the same end of the text it
was on: after the off code if it was at the end of the block, after the on
code if it was at the start.

\Topic{The Font Menu With Block On}
With Block on, the Font menu offers only Size and Appearance.  Normal is not
offered.

\SeeAlso{Appearance; Block; Bold; Size; Underline}
\end{Feature}

\begin{Feature}{Block, Delete}
You can delete a whole block of text at once.

To delete a block,
\begin{WPSteps}
\item Block the text (see \InRef{Block}).
\item Press \K{Delete} or \K{Backspace}.  The status line asks
\Scr{Delete Block? No (Yes)}.
\item Type \Typed{y} to delete the block.  Any other key leaves the block as
it is, with Block still on.
\end{WPSteps}
\Documents{wpkey}{Block:Delete}
\Documents{wpkey}{Block:Backspace}

\Notes
\Topic{Codes in the Block}
Text outside the block keeps the appearance and size it had.  If the block
held the start or the end of an attribute, InitechWord puts the codes needed
for that where the block was.

\Topic{Delete With Move}
You can also delete a block with \FK{Move}{Ctrl-F4}: select \Ch{Block}{1},
then \Ch{Delete}{3} (see \InRef{Move, Block}).  Move deletes without asking.

\SeeAlso{Block; Delete; Move, Block}
\end{Feature}

\begin{Feature}{Bold}
Bold makes text stand out by printing it darker.  On the editing screen
bold text is drawn in bold.

To make text bold as you type it,
\begin{WPSteps}
\item Press \FK{Bold}{F6}.  InitechWord puts the codes \Code{[BOLD]} and
\Code{[bold]} in the document, with the cursor between them.
\item Type the text.
\item Press \FK{Bold}{F6} again.  The cursor moves past \Code{[bold]}, so
that what you type next is not bold.
\end{WPSteps}
\Documents{wpkey}{F6}
\begin{WPNote}To make text bold after you have typed it, block it and press
Bold (see \InRef{Block, Attributes}).\end{WPNote}

\BeginReplay{ref-bold}
For example, type \TypeRun{Please read the }(with the space), press
\Press{Bold}{F6}, type \TypeRun{new}, press \Press{Bold}{F6} again and type
\TypeRun{ policy.}  The screen shows:
\WordScreen[rows=1]{ref-bold}{A bold word typed with Bold (F6)}
\EndReplay

\Notes
\Topic{Bold From the Font Menu}
Bold is also Bold (1) on the Appearance menu of the Font key: \KC{Ctrl-F8}{2,1}
does what \K{F6} does.

\Topic{Bold in Bold Text}
Pressing Bold with the cursor in bold text ends the bold at the cursor: the
\Code{[bold]} code is moved to the cursor and the cursor goes past it.

\Topic{Removing Bold}
Delete the \Code{[BOLD]} or the \Code{[bold]} code; the other one goes with
it (see \InRef{Codes}).

\SeeAlso{Appearance; Block, Attributes; Codes; Underline}
\end{Feature}
